\documentclass[10pt,a4paper]{amsart}
\usepackage{amsmath,amssymb,amsthm,graphicx}
% Minimal retained preamble for lane-s5's SIGMA 2025-022 cards, taken from
% sources/sigma-2025-022-mode-stability-of-hermitian-instantons/pinned-original-source.tex
% (doi:10.3842/SIGMA.2025.022, SIGMA 21 (2025), article 022, CC BY 4.0).
%
% Retained from the source preamble: the equation numbering and the theorem declarations under
% the shared Theorem counter, verbatim.
%
% Omitted and why: \pdfoutput=1 and the \ifpdf / \documentclass[pdftex]{sigma} block (pdfTeX
% primitive and the journal class, which texlive does not carry; per the lane-x1 precedent the
% assembler supplies the baseline class), the SIGMA front-matter commands and the paper's own
% \arXivNumber/\PaperNumber definitions, none of which the retained spans use.
\usepackage{amsmath,amssymb,amsthm,mathrsfs,mathtools}
\usepackage{hyperref}% the author bibliography entries use \href (assembler lesson 8.8.157)

\numberwithin{equation}{section}

\newtheorem{Theorem}{Theorem}[section]
\newtheorem*{Theorem*}{Theorem}
\newtheorem{Corollary}[Theorem]{Corollary}
\newtheorem{Lemma}[Theorem]{Lemma}
\newtheorem{Proposition}[Theorem]{Proposition}
\newtheorem{Conjecture}[Theorem]{Conjecture}
{ \theoremstyle{definition}
\newtheorem{Definition}[Theorem]{Definition}
\newtheorem{Note}[Theorem]{Note}
\newtheorem{Example}[Theorem]{Example}
\newtheorem{Remark}[Theorem]{Remark} }
\title{The conformal operator identity behind the Teukolsky equation for Einstein--Hermitian perturbations}
\author{Andersson, Lars, Araneda, Bernardo, Dahl, Mattias}
\date{doi:10.3842/SIGMA.2025.022; CC BY 4.0}
\begin{document}
\maketitle
\setcounter{section}{2}
\setcounter{subsection}{1}
\setcounter{Theorem}{0}
\setcounter{equation}{0}
\noindent\textit{Editorial note.} All delivered mathematical text is the author's own, copied from the pinned SIGMA article source (source\_id doi:10.3842/SIGMA.2025.022); no statement or proof step was rewritten, reordered or removed.\par\medskip
\noindent\textit{Editorial note.} The author's wording is delivered verbatim, including any typographical peculiarity of the source; nothing was silently corrected (CONTROL 8.8.124).\par\medskip
\noindent\textit{Editorial note.} The author proves this lemma in a named proof placed after the following Remark; the proof span is that named proof (CONTROL 8.8.157: a far-placed named proof is delivered as the single proof span).\par\medskip
\noindent\textit{Editorial note.} The spans cite PR:I, PR:II and araneda2018conformal; the card carries a card-local bibliography.\par\medskip
\noindent\textit{Editorial note.} The card carries a minimal card-local preamble instead of the article's own header (which uses the pdfTeX-only \textbackslash{}pdfoutput primitive and the journal class sigma.cls); the retained lines are copied verbatim from the pinned source header (CONTROL 8.8.165).\par\medskip
\noindent\textit{Editorial note.} Referenced displays of the same paper are delivered as context spans so that the reader sees the equations the proof compares against (lane-s5 paper convention).\par\medskip
\noindent\textit{Editorial note.} The preamble span points at the card-local preamble; the counted spans are the author's own environments in the pinned source.\par\medskip
\noindent\textit{Allowed standard prerequisites (editorial).}\begin{itemize}
\item GHP formalism and the weighted operators acting on scalar fields. Same paper, Section 2, in force in the statement. Fixes the notation in which the identity is written.
\item The Teukolsky operator and the linearised Einstein equation. Same paper, the displayed equations preceding the lemma (statement span references them). Supplies the two sides of the identity.
\end{itemize}\medskip
\begin{align}\label{eq:EinsteinOp}
E_{ab}(g) = R_{ab}(g) - \frac{S(g)}{4}g_{ab},
\end{align}

\begin{align} \label{eq:LinMetric}
\frac{{\rm d}}{{\rm d}s} g(s)_{ab} |_{s=0} = h_{ab},
\end{align}

\begin{align} \label{eq:LinPsi0}
\vartheta \Psi_0 = \frac{{\rm d}}{{\rm d}s} \Psi_0[g(s)_{ab}] |_{s=0},
\end{align}

\subsection{Spinors and complex structures} \label{sec:spinors}

\subsection{Conformally invariant GHP connections} \label{sec:CGHP}

\begin{align}
&f_{a} = -\frac{1}{2}J^{c}{}_{b}\nabla_{c}J^{b}{}_{a}, \label{LeeForm} \\
&P_{a} = \omega_{a} -\frac{1}{2} {\rm i} J^{b}{}_{a}f_{b}, \label{1formP} \\
&Q_{ab}{}^{c} = f_{a}\delta^{c}_{b} + f_{b}\delta^{c}_{a} - f^{c}g_{ab}. \nonumber
\end{align}

\begin{align}\label{SFR}
o^{A}o^{B}\nabla_{AA'}o_{B} = 0 \qquad \Leftrightarrow \qquad \kappa=\sigma=0,
\end{align}

\begin{align}
\vartheta E_{ab} &{}= \vartheta R_{ab} - \frac{1}{4}\vartheta S g_{ab} - \frac{1}{4}S(g)h_{ab} \nonumber\\
&{}= \frac{1}{2}\Delta h_{ab} - \frac{1}{2}\nabla_{a}\nabla_{b}\bigl(g^{cd}h_{cd}\bigr) + \frac{1}{2}\nabla^{c}\nabla_{a}h_{bc} + \frac{1}{2}\nabla^{c} \nabla_{b}h_{ac} \nonumber\\
&\quad{} - \frac{1}{4}g_{ab}\bigl[ \nabla^{c}\nabla^{d}h_{cd} + \Delta\bigl(g^{cd}h_{cd}\bigr) - h^{cd}R_{cd} \bigr]
-\frac{1}{4}S(g) h_{ab},\label{LinEinsteinOp}
\end{align}

\begin{align}\label{TeukOp}
 {\rm L}=g^{ab}\mathcal{C}_{a}\mathcal{C}_{b}-18\Psi_{2}
\end{align}

\begin{Lemma}\label{lem:OpId}
Let $(M,g_{ab})$ be Einstein--Hermitian, with $($possibly vanishing$)$ cosmological constant $\lambda$. Let $h_{ab}$ be an arbitrary metric perturbation \eqref{eq:LinMetric}, $\vartheta{E}_{ab}$ the linearized Einstein operator~\eqref{LinEinsteinOp}, and $\vartheta\Psi_0$ the linearized Weyl scalar~\eqref{eq:LinPsi0}. Furthermore, let $f_{a}$ be the Lee form \eqref{LeeForm}, and $Q^{abcd}$ the tensor field
\begin{align}\label{tensorQ}
Q^{abcd} = o^{A}o^{B}o^{C}o^{D}\epsilon^{A'B'}\epsilon^{C'D'}.
\end{align}
Then
\begin{align}\label{eq:OpIdentity}
-\mathring{\Omega}^{-1}Q^{acbd}(\nabla_{a}-4f_{a})\nabla_{d}\vartheta E_{bc} = {\rm L}\bigl[\mathring{\Omega}^{-1}\vartheta\Psi_0\bigr],
\end{align}
where ${\rm L}$ is the operator \eqref{TeukOp} and $\mathring{\Omega}$ is an auxiliary constant conformal factor, that is a scalar field with weights $w=1$, $p=0$ and $\nabla_{a}\mathring\Omega=0$.
\end{Lemma}

\begin{proof}[Proof of Lemma \ref{lem:OpId}]
The strategy is to consider identities for an arbitrary Riemannian manifold $(M,g_{ab})$, and then to linearize around an Einstein--Hermitian metric. Consider then an arbitrary $(M,g_{ab})$, with Levi-Civita connection $\nabla_{a}$ and unprimed Weyl curvature spinor $\Psi_{ABCD}$. Let $J^{a}{}_{b}$ be a (locally defined) compatible almost-complex structure, and let $o^{A}$ be the associated spinor field as in Section \ref{sec:spinors}. We can then define conformally and GHP weighted fields as in Section \ref{sec:CGHP}, together with the connection $\mathcal{C}_{a}$ on the corresponding bundles.

Let $\mathring\Omega$ be an arbitrary constant conformal factor, that is a scalar field with weights~${w=1}$, ${p=0}$ and $\nabla_{a}\mathring\Omega=0$, and define
\begin{align}\label{GravSpin2Field}
 \varphi_{ABCD} := \mathring\Omega^{-1}\Psi_{ABCD}.
\end{align}
This object has weights $w=-1$, $p=0$, and it is essentially the ``gravitational spin 2 field'' of Penrose and Rindler \cite[equation~(9.6.40)]{PR:II}. We have
\begin{align*}
 \mathcal{C}_{AA'}\mathcal{C}^{A'E}\varphi_{BCDE} &=
(\nabla_{AA'}-4f_{AA'})\nabla^{A'E}\varphi_{BCDE} \\
 &= \mathring\Omega^{-1}(\nabla_{AA'}-4f_{AA'})\nabla^{A'E}\Psi_{BCDE} \\
 &= \mathring\Omega^{-1}(\nabla_{AA'}-4f_{AA'})\nabla_{(B}^{B'}\Phi_{CD)B'}{}^{A'},
\end{align*}
where in the second line we used the definition \eqref{GravSpin2Field} and the fact that $\mathring\Omega$ is constant, and in the third line we used the spinor form of the Bianchi identities \cite[equation~(4.10.7)]{PR:I}
adapted to Riemann signature (recall Section \ref{sec:spinors}).
Here $\Phi_{ABA'B'}$ is the trace-free Ricci spinor. Contracting with $o^Ao^Bo^Co^D$:
\begin{align}\label{AuxId}
\nonumber o^Ao^Bo^Co^D\mathcal{C}_{AA'}\mathcal{C}^{A'E}\varphi_{BCDE} &{} = \mathring\Omega^{-1}o^Ao^Bo^Co^D (\nabla_{AA'}-4f_{AA'})\nabla_{B}^{B'}\Phi_{CDB'}{}^{A'} \\
\nonumber &{} = \mathring\Omega^{-1}o^Ao^Co^Bo^D\epsilon^{A'C'}\epsilon^{B'D'} (\nabla_{AA'}-4f_{AA'})\nabla_{DD'}\Phi_{BCB'C'}\\
&{} = -\frac{1}{2}\mathring\Omega^{-1}Q^{acbd}(\nabla_{a}-4f_{a})\nabla_{d}E_{bc},
\end{align}
where in the last line we used the definition \eqref{tensorQ} and the identity \smash{$\Phi_{bc}=-\frac{1}{2}E_{bc}$} (see \cite[equa\-tion~(4.6.25)]{PR:I} and recall \eqref{eq:EinsteinOp}).

It remains to find a convenient expression for the first line in \eqref{AuxId}. This can be done using \cite[equations~(3.6) and (3.15)]{araneda2018conformal}. We have
\begin{align}\label{AuxId2}
o^Ao^Bo^Co^D\mathcal{C}_{AA'}\mathcal{C}^{A'E}\varphi_{BCDE} =
\frac{1}{2} \bigl(g^{ab}\mathcal{C}_{a}\mathcal{C}_{b}-18\Psi_{2}\bigr)\varphi_{0} + B,
\end{align}
where $\varphi_{0}= o^Ao^Bo^Co^D\varphi_{ABCD}=\mathring{\Omega}^{-1}\Psi_0$ and $B$ is a term which couples the GHP quantities $\kappa$, $\sigma$, $\Psi_0$, $\Psi_1$ quadratically. Equating the last line of~\eqref{AuxId} to the right-hand side of~\eqref{AuxId2}, and taking a~linearization around a metric which satisfies $E_{ab}|_{s=0}=0$, ${\Psi_{0}|_{s=0}=\Psi_{1}|_{s=0}=\kappa|_{s=0}=\sigma|_{s=0}=0}$, that is, an Einstein--Hermitian metric, see \eqref{SFR}, the result \eqref{eq:OpIdentity} follows.
\end{proof}

% Card-local bibliography so that the citations carried by this card's spans resolve.
% Every \bibitem below is copied character for character from the pinned source's own
% thebibliography (sources/sigma-2025-022-mode-stability-of-hermitian-instantons/pinned-original-source.tex).
% Only the keys cited by this card are retained.
\begin{thebibliography}{99}

\bibitem{araneda2018conformal}
Araneda B., Conformal invariance, complex structures and the {T}eukolsky
 connection, \href{https://doi.org/10.1088/1361-6382/aad13b}{\textit{Classical Quantum Gravity}} \textbf{35} (2018), 175001,
 23~pages, \href{https://arxiv.org/abs/1805.11600}{arXiv:1805.11600}.

\bibitem{PR:I}
Penrose R., Rindler W., Spinors and space-time. {V}ol.~1. {T}wo-spinor calculus
 and relativistic fields, \textit{Cambridge Monogr. Math. Phys.}, \href{https://doi.org/10.1017/CBO9780511564048}{Cambridge University
 Press}, Cambridge, 1984.

\bibitem{PR:II}
Penrose R., Rindler W., Spinors and space-time. {V}ol.~2. Spinor and twistor
 methods in space-time geometry, \textit{Cambridge Monogr. Math. Phys.}, \href{https://doi.org/10.1017/CBO9780511524486}{Cambridge
 University Press}, Cambridge, 1986.

\end{thebibliography}
\end{document}
